\documentclass[10pt]{article}
\usepackage{amsmath}
\usepackage{amssymb}
\usepackage{comment}
\usepackage{fancyhdr}
\usepackage{hyperref}
\usepackage{listings}
\usepackage{longdivision}
\usepackage{booktabs}
\usepackage{polynom}
\usepackage{graphicx}
\usepackage{tikz}
\usepackage{float}
\usepackage{cancel}
\usepackage{xlop}
\usepackage{xcolor} % for colors


\setlength{\parindent}{0pt}
\setlength{\parskip}{0.5em}

% Add license
\usepackage{ccicons}
\usepackage[
    type={CC},
    modifier={by-nc-sa},
    version={4.0},
]{doclicense}

\pagestyle{fancy}
\fancyhf{} % clear all header and footer fields

% Footer center
\fancyfoot[C]{\small\textcolor{gray}{
GitHub repository: \href{https://github.com/deepak-venkatesh/gelfand-method-coordinates}{https://github.com/deepak-venkatesh/gelfand-method-coordinates}
}}


% Page number (optional)
\fancyfoot[R]{\thepage}

\renewcommand{\headrulewidth}{0pt}
\renewcommand{\footrulewidth}{0.4pt}


\title {An Unofficial Compendium of Worked Solutions to the Problems in \textit{The Method of Coordinates}\\
[0.5em]\large by I.M. Gelfand, E. G. Glagoleva \& A. A. Kirillov}

\author{Deepak Venkatesh}
\date{\today}


\begin{document}
\noindent
\maketitle
\begin{abstract}
\noindent
This book is an independent and unofficial collection of worked solutions
to problems inspired by the book \textit{The Method of Coordinates} by I.\,M.~Gelfand, E.\,G.~Glagoleva and A.\,A.~Kirillov,
first published in 1967. The original text presents a wide range of problems
in elementary algebra, some accompanied by solutions and others left as exercises
for the reader.\\

\noindent
The aim of this book is to provide clear, logically consistent solutions
to all problems addressed, written in a style accessible to motivated students
and self learners. Proofs have been verified primarily by hand, and in some cases
supported by small computational checks written in Scheme (a dialect of Lisp) or in Python.
Large language models were used only as auxiliary tools, primarily for generating
figures in \LaTeX\ and for reviewing exposition.\\

\noindent
This work is not affiliated with, endorsed by, or approved by the original authors
or publisher. Any errors, omissions, or inconsistencies are entirely my own.
Issues and corrections may be reported at
\href{https://github.com/deepak-venkatesh/gelfand-method-coordinates}{this repository}.
\end{abstract}

\clearpage
\thispagestyle{empty}

\vspace*{2cm}

\begin{center}
  {\Large \textbf{Copyright and Disclaimer}}
\end{center}

\vspace{1cm}

\noindent
\textbf{Unofficial companion.}

\vspace{0.4cm}

\noindent
This document is an \textbf{independent, unofficial collection of worked solutions}
to problems from \textit{The Method of Coordinates} by I.\,M.~Gelfand , E.\,G.~Glagoleva and A.\,A.~Kirillov (first published 1967).

\vspace{0.4cm}

\noindent
This work is \textbf{not affiliated with, endorsed by, or approved by} the original
authors or the publisher. The names of the authors and the title of the book are used
only for identification and reference.

\vspace{0.4cm}

\noindent
The original problem statements are not reproduced here. All exposition and solutions
are written independently by the author of this document.

\vspace{1.5cm}

\noindent
\textbf{Copyright \textcopyright\ \the\year\ Deepak Venkatesh}

\vspace{1.5cm}

\noindent
License: \doclicenseThis

\vfill

\noindent
Project repository:\\
\href{https://github.com/deepak-venkatesh/gelfand-method-coordinates}{https://github.com/deepak-venkatesh/gelfand-method-coordinates}

\clearpage

\begin{comment}
  \vspace{27em}
  \doclicenseThis
\end{comment}

\newpage
\thispagestyle{empty}

\vspace*{\fill}
\begin{center}
\textit{For\\[0.5em]
My Daughter and Son}
\end{center}
\vspace*{\fill}
\newpage
\tableofcontents
\newpage % start sections on a fresh page

\textbf{Working Notes}\\\\

\emph{Challenging Problems for Middle Schoolers:}\\

Problem Numbers:  \\\\

\emph{Typographical errors in:}\\

Problem Numbers: 

\newpage

\part{Part 1}

\section{Chapter 1: The Coordinates of Points on a Line}

%%%%%%%%%% 1
\subsection{The Number Axis}

\textbf{Exercise 1}

\textbf{(a)}

The points are marked below. Note that $\dfrac{13}{3}$ lies somewhere between 4 and 5.

\begin{center}
\begin{tikzpicture}[scale=1.2]
    % Number line
    \draw[<->, thick] (-3,0) -- (5.5,0);

    % Points
    \fill (-2,0) circle (1pt);
    \fill (0,0) circle (1pt);
    \fill (13/3,0) circle (1pt);

    % Labels
    \node[above] at (-2,0) {$A$};
    \node[below] at (-2,0) {$-2$};

    \node[above] at (0,0) {$K$};
    \node[below] at (0,0) {$0$};

    \node[above] at (13/3,0) {$B$};
    \node[below] at (13/3,0) {$\frac{13}{3}$};
\end{tikzpicture}
\end{center}

\textbf{(b)}

The two points are 3 units to the right and 3 units to the left of $M$ as shown in the line below.

\begin{center}
\begin{tikzpicture}[scale=1.2]
    % Number line
    \draw[<->, thick] (-3,0) -- (6,0);

    % Points
    \fill (-1,0) circle (1pt);
    \fill (2,0) circle (1pt);
    \fill (5,0) circle (1pt);

    % Labels
    \node[above] at (2,0) {$M$};
    \node[below] at (2,0) {$2$};

    \node[above] at (-1,0) {$A$};
    \node[below] at (-1,0) {$-1$};

    \node[above] at (5,0) {$B$};
    \node[below] at (5,0) {$5$};
\end{tikzpicture}
\end{center}

\textbf{Exercise 2}

\textbf{(a)}

Going by the convention that the right side of the number line indicates larger nummber, in this 
exercise $a > b$. \\

\textbf{(b)}

On the right are the larger numbers
\begin{itemize}
  \item $A(-3)$
  \item $B(4)$
  \item $B(4)$
  \item $A(3)$
\end{itemize}


\textbf{Exercise 3}

It is easy to say that $A(a)$ is on the right of $B(-a)$ but this assumes that $a$ is a positive number.
So this could be true but the other case is that if $a$ were a negative number then $A(a)$ would be a negative number 
and $B(-a)$ would be a positive number, thus $B(-a)$ would be on the right.

The third case is that $a = 0$ then both $A$ and $B$ would be 0 and they would be at the same point on the number line.\\


\textbf{Exercise 4}

\textbf{(a)}

If $x = 0$, both $M(x)$ and $N(2x)$ will be at the same point on the number line.\\

If $x > 0$, $N(2x)$ will be on the right of $M(x)$ on the number line.\\

If $x < 0$, $M(x)$ will be on the right of $N(2x)$ on the number line.\\


\textbf{(b)}

In this problem no matter whether $c$ is 0, positive or negative $B(c + 2)$ will always be on the right side 
of $A(c)$.\\

\textbf{(c)}

In this problem we have 3 scenarios.\\

If $a = 0$ then $A(x)$ and $B(x - a)$ are the at the same point in the number line.\\

If $a > 0$ then $A(x)$ will be on the right of $B(x - a)$.\\

If $a < 0$ then $B(x - a)$ will be on the right side of $A(x)$ since the negative number multiplied by $-$ will give a positive number. \\

\textbf{(d)}

If $x = 0$ then both $A(x)$ and $B(x^2)$ would be at the same point in the number line.\\

But in any other case whether $x$ is positive or negative $x^2$ will always be greater (and positive) that $x$. Thus 
$B(x^2)$ will always be on the right in these two cases.\\


\textbf{Exercise 5}

The space or distance between the two points is 12 (-5 to 0 and 0 to 7). The middle point will be 6 units from -5 on the right which is at point 1.\\

\begin{center}
\begin{tikzpicture}[scale=1.0]
    % Number line
    \draw[<->, thick] (-6,0) -- (8,0);

    % Points
    \fill (-5,0) circle (1pt);
    \fill (1,0) circle (1pt);
    \fill (7,0) circle (1pt);

    % Labels
    \node[above] at (-5,0) {$A$};
    \node[below] at (-5,0) {$-5$};

    \node[above] at (1,0) {$\text{middle}$};
    \node[below] at (1,0) {$1$};

    \node[above] at (7,0) {$B$};
    \node[below] at (7,0) {$7$};
\end{tikzpicture}
\end{center}


\textbf{Exercise 6}

\textbf{(a)}

Whole numbers will include number 0.

\begin{center}
\begin{tikzpicture}[scale=1.0]
    % Number line
    \draw[<->, thick] (-6,0) -- (9,0);

    % Mark and label whole numbers
    \foreach \x in {0,1,2,3,4,5,6,7}
    {
        \fill[red] (\x,0) circle (1.5pt);
        \node[below] at (\x,0) {$\x$};
    }

    % Dots after 7
    \node[below] at (8,0) {$\cdots$};
\end{tikzpicture}
\end{center}

\textbf{(b)}

Positive numbers will not include number 0.

\begin{center}
\begin{tikzpicture}[scale=1.0]
    % Number line
    \draw[<->, thick] (-6,0) -- (9,0);

    % Mark and label whole numbers
    \foreach \x in {1,2,3,4,5,6,7}
    {
        \fill[red] (\x,0) circle (1.5pt);
        \node[below] at (\x,0) {$\x$};
    }

    % Dots after 7
    \node[below] at (8,0) {$\cdots$};
\end{tikzpicture}
\end{center}


\textbf{Exercise 7}

\textbf{(a)}

For $x < 2$\\

\begin{center}
\begin{tikzpicture}[scale=1.0]

    % Number line
    \draw[<->, thick] (-2,0) -- (5,0);

    % Vertical connector
    \draw[blue, thick] (2,0) -- (2,0.7);

    % Solution arrow pointing left
    \draw[blue, ->, very thick] (2,0.7) -- (-1.5,0.7);

    % Open circle
    \draw[blue, thick, fill=white] (2,0.7) circle (3pt);

    % Label
    \node[below] at (2,0) {$2$};

\end{tikzpicture}
\end{center}


\textbf{(b)}

For $x \ge 5$\\

\begin{center}
\begin{tikzpicture}[scale=1.0]

    % Number line
    \draw[<->, thick] (2,0) -- (9,0);

    % Vertical connector
    \draw[blue, thick] (5,0) -- (5,0.7);

    % Solution arrow pointing right
    \draw[blue, ->, very thick] (5,0.7) -- (8.5,0.7);

    % Filled circle
    \fill[blue] (5,0.7) circle (3pt);

    % Label
    \node[below] at (5,0) {$5$};

\end{tikzpicture}
\end{center}


\textbf{(c)}

For $2 < x < 5$\\

\begin{center}
\begin{tikzpicture}[scale=1.0]

    % Number line
    \draw[<->, thick] (0,0) -- (7,0);

    % Vertical connectors
    \draw[blue, thick] (2,0) -- (2,0.7);
    \draw[blue, thick] (5,0) -- (5,0.7);

    % Solution interval
    \draw[blue, very thick] (2,0.7) -- (5,0.7);

    % Open circles
    \draw[blue, thick, fill=white] (2,0.7) circle (3pt);
    \draw[blue, thick, fill=white] (5,0.7) circle (3pt);

    % Labels
    \node[below] at (2,0) {$2$};
    \node[below] at (5,0) {$5$};

\end{tikzpicture}
\end{center}


\textbf{(d)}

For $-3\frac{1}{4} \le x \le 0$\\

\begin{center}
\begin{tikzpicture}[scale=1.0]

    % Number line
    \draw[<->, thick] (-5,0) -- (2,0);

    % Vertical connectors
    \draw[blue, thick] (-3.25,0) -- (-3.25,0.7);
    \draw[blue, thick] (0,0) -- (0,0.7);

    % Solution interval
    \draw[blue, very thick] (-3.25,0.7) -- (0,0.7);

    % Filled circles
    \fill[blue] (-3.25,0.7) circle (3pt);
    \fill[blue] (0,0.7) circle (3pt);

    % Labels
    \node[below] at (-3.25,0) {$-3\frac{1}{4}$};
    \node[below] at (0,0) {$0$};

\end{tikzpicture}
\end{center}

%%%%%%%%%%%%%%%%%%%%%%%%%%%%%%%%%%%%%%%%%
% Chapter End
%%%%%%%%%%%%%%%%%%%%%%%%%%%%%%%%%%%%%%%%%

%%%%%%%%%% 2
\subsection{The Absolute Value of a Number}

\textbf{Exercise 1}

There will be 3 cases:\\

Case 1: $x = 0$. In this we are dividing by 0 so the division is undefined.\\

Case 2: $x < 0$. The expression is $\dfrac{x}{-x} = -1$. It is -1.\\

Case 3: $x > 0$. Here it is simply 1\\

\textbf{Exercise 2}

\textbf{(a)}
The square will always be positive. Hence the term $a^2$ with or without modulus will always be positive.
We can write it as.
\[
|a^2| = a^2
\] 

\textbf{(b)}
The term will always be positive since $a > b$. 
\[
|a - b| = a - b
\]

\textbf{(c)}
In this case the term will be negative since $a < b$.
\[
|a - b| = - (a - b) = b - a
\]

\textbf{(d)}
$a$ is negative so negative of it will be positive, therefore we will get the below.
\[
|- a| = a
\]


\textbf{Exercise 3}

We know that the term $x - 3$ will be non negative. So we can say that $x \ge 3$.




%%%%%%%%%% 3
\subsection{The Distance Between Two Points}



\section{Chapter 2: The Coordinates of Points in the Plane}

%%%%%%%%%% 4
\subsection{The Coordinate Plane}

%%%%%%%%%% 5
\subsection{Relations Connecting Coordinates}

%%%%%%%%%% 6
\subsection{The Distance Between Two Points}

%%%%%%%%%% 7
\subsection{Defining Figures}

%%%%%%%%%% 8
\subsection{We Begin to Solve Problems}

%%%%%%%%%% 9
\subsection{Other System of Coordinates}



\section{Chapter 3: The Coordinates of a Point in Space}

%%%%%%%%%% 10
\subsection{Coordinate Axes and Planes}

%%%%%%%%%% 11
\subsection{Defining Figures in Space}



\part{Part 2}

\section{Chapter 1: Introduction}

%%%%%%%%%% 1
\subsection{Some General Considerations}

%%%%%%%%%% 2
\subsection{Geometry as an Aid in Calculation}

%%%%%%%%%% 3
\subsection{The Need for Introducing Four-Dimensional Space}

%%%%%%%%%% 4
\subsection{The Peculiarities of Four-Dimensional Space}

%%%%%%%%%% 5
\subsection{Some Physics}


\section{Chapter 2: Four-Dimensional Space}

%%%%%%%%%% 6
\subsection{Coordinate Axes and Planes}

%%%%%%%%%% 7
\subsection{Some Problems}


\section{Chapter 3: The Four-Dimensional Cube}

%%%%%%%%%% 8
\subsection{The Definition of the Sphere and the Cube}

%%%%%%%%%% 9
\subsection{The Structure of the Four-Dimensional Cube}

%%%%%%%%%% 10
\subsection{Problems on the Cube}

 
\end{document}